\section{Conclusion}
\label{sec:conclusion}

Input data manipulation is a genuine gaming vector for algorithmic
redistricting: an adversary with access to the data pipeline can modify
population counts, shapefile boundaries, or demographic tables before
the algorithm processes them.
However, the vector is both bounded and defended.

Bounded: even a combined attack across all three data sources produces
at most $\Delta D \leq 0.4$ Democratic seats nationally, well within
the B.7 seed-variance noise floor.
The constraint is geometric---competitive-district geography cannot
absorb more than a small fraction of the adversary's manipulation
before the districts are non-competitive anyway.

Defended: \bisect{}'s three-level SHA-256 audit chain (download hash,
adjacency hash, plan hash) makes undetected manipulation computationally
infeasible.
Any modification that exceeds the $\pm 1\%$ noise threshold produces
a hash mismatch detectable by any party who independently downloads
and verifies the Census source files.

The disclosed TLS limitation (no certificate pinning) is a real but
bounded risk.
It requires infrastructure-level compromise of the redistricting
computing environment, which is a broader security failure not specific
to the data manipulation gaming vector.
The recommended mitigation (multi-source hash verification) eliminates
this risk for litigation-grade deployments.

The combined finding of R.1 and R.2 for the Track R program: neither
parameter gaming nor input data manipulation can produce a meaningful
partisan outcome undetected.
The remaining plan-affecting gaming vector (V3: geographic definition)
is analysed in R.3; the audit chain formalization as a legal defense
is developed in R.4.
